\documentclass[ekobook.tex]{subfiles}

\begin{document}\setWPP
\setlist[itemize]{
  itemsep=0pt,
  parsep=0pt,      % space between paragraph lines inside one item
  topsep=0.5ex,    % space before/after the whole list (usually keep a little)
  partopsep=0pt
}
\lsection{Strategie prodeje a cenová politika}
\qaw{Cílený marketing}{%
Firma si rozdělí trh (a tedy zákazníky) \textbl{na segmenty} dle určitých kritérií.
Geopolitická situace, gemografie či dle sociálněpsychologické situace.
Ke každému segmentu má mírně \texthl{odlišný přístup} a \texthl{marketingovou strategii}.
}
\qaw{Druhy segmentace}{%
\begin{itemize}[label={\color{WPP}\faSlackHash}]
\item \textbl{geografické členění}
-- jinak nakupují lidé ve Grónsku a jinak v Itálii
\item \textbl{gemografické členění}
-- jinak nakupují puberťáci, jinak padesátníci; jinak muži, jinak ženy; jinak úředníci, jinak vědci, jinak učitelé, jinak pracovníci u pásu; poptávka se liší i dle národnosti (například kvůli jiným chutěm a jiné národní kuchyni)
\item \textbl{sociálněpsychologické členění}
-- příjem, životní styl a další
\end{itemize}
}
\qaw{CMR}{%
\textbl{Customer Managed Relationship}, zákazníkem řízený vztah, systém a strategie pro sběr, anakýzu a využití dat o zákazníkovi.\\
\vspace{1pt}\\
\textbl{získání individuálních dat}
-- zákazník musí podepsat GDPR, jinak je sběr dat ilegální (například souhlas s podmínkami na e--shopu)\\
\textbl{analýza dat}\\
\textbl{individuálně přizpůsobená nabídka}
-- například cílená reklama dle toho, co si zákazník prohlížel\\
\textbl{realizace dodávky}
-- odhadneme potřeby služby zákazníka (doprava, instalace, servis)\\
\textbl{křížový prodej}
-- zákazníkovi časem nabídneme související produkty\\
\textbl{up-sell}
-- cíleně nabídneme novou generaci produktů lidem, co si koupili ze staré generace (například Apple)
}
\qaw{Co umí sledovač prodeje}{%
Sledovač sleduje, kdy lidé \textbl{nejčastěji nakupují} a v tyto doby \textbl{zvyšuje cenu}.
E-shop zvýší cenu zákazníkům, \texthl{co pravidelně nakupují} (plánuje to například Sony u svého Playstationu, věrní hráči budou potrestáni vyššími cenami, protože pravidelně nakupují).
Aerolinky zvýší cenu letenek na lety, kterými \textbl{častěji lidé létají} (případně zvýší cenu v hodinu, kdy lidé nejčastěji objednávají, u některých firem se vyplatí nakupovat po půlnoci).
V kamenných obchodech již existují cenovky jako digitální displeje, které jsou schopny zobrazovat změněné ceny během dne; ceny fluktuují dle fluktuace zákazníků.
}
\qaw{Co ve strategiích prodeje využívají firmy}{%
\textbl{Umělá inteligence} je schopna rychle a efektivně odhadovat chování zákazníků na základě analýzy dat, může tak \texthl{efektivně měnit ceny} a \texthl{podbízet další nabídky}, u kterých je téměř jistota, že si je zákazník koupí.\\
\textbl{Průběžné sledování komunikace}, kdy e-shopy sledují \texthl{klikání zákazníků}, čas strávený \texthl{sledováním reklamy}, co si zákazník \texthl{dal do košíku, ale nekoupil}; na základě všech těchto dat pak mění ceny a vytvářejí slevové akce.\\
\textbl{Personalizace nabídky} reklama na webech, e-mailem i doporučená nabídka na stránkách e-shopů, \texthl{se postupem času mění na míru zákazníkovi}, za účelem maximalizace prodeje.
}
\qaw{Cena -- způsoby stanovení ceny (získat nové trhy, jako konkurence, politika vyšších cen)}{%
\textbl{chci se dostat na nové trhy}
-- stanovím cenu na nižší, než je tržní, tím nalákám zákazníka od konkurence\\
\textbl{prodávání za ceny jako konkurence}
-- cenu stanovím na tržní, konkuruji pouze kvalitou výrobků\\
\textbl{politika vyšších cen}
-- cenu stanovím vysoko nad tržní, tím dávám najevo, že se jedná o kvalitní zboží, a cílím pouze na menší skupinu zákazníků, co si mohou zboží dovolit
}
\qaw{Metody stanovení ceny}{%
podle \texthl{nákladů}; kdy propočítáme náklady na výrobu a cenu přizpůsobíme tomu\\
podle \texthl{orientace na konkurenci}; sledujeme tržní cenu a dle naší politiku stanovíme nižší, stejnou nebo vyšší cenu\\
podle \texthl{hodnoty vnímané zákazníkem}; provedeme průzkum trhu a stanovíme cenu dle toho, jakou cenu výrobku přisuzuje zákazník (praktická jistota prodeje)\\
podle \texthl{orientace na poptávku}; čím je cena nižší, tím větší poptávka
}
\qaw{Strategie pronikání a sbírání smetany}{%
\texthl{strategie pronikání}; rychlé proniknutí na trh, kde podobné zboží již existuje a my přetáhneme zákazníky svojí nízkou cenou\\
\texthl{strategie sbírání smetany}; výrobek je odlišný od konkurence, vysokou zaváděci cenou naznačíme jeho morální hodnotu, cílíme na inovačně zaměřené zákazníky (Mercedess, Apple, ...)
}
\wpage
\end{document}